\documentclass[11pt]{article}
\usepackage{amsmath,textcomp,amssymb,geometry,graphicx,enumerate}

\def\Name{Vijay Hans}  % Your name
\def\SID{3041285081}  % Your student ID number
\def\PreLab{1} % Number of Homework
\def\Session{Spring 2026}
\def\Cls{PHYS 5BL}
\title{\Cls--Spring 2026 --- Pre-Lab \PreLab \ Response}
\author{\Name, SID \SID}
\markboth{\Cls -- \Session\  Pre-Lab \PreLab\ \Name}{\Cls -- \Session\ Pre-Lab \PreLab\ \Name}
\pagestyle{myheadings}
\date{\today}

\newenvironment{qparts}{\begin{enumerate}[{(}a{)}]}{\end{enumerate}}
\def\endproofmark{$\Box$}
\newenvironment{proof}{\par{\bf Proof}:}{\endproofmark\smallskip}

\textheight=9in
\textwidth=6.5in
\topmargin=-.75in
\oddsidemargin=0.25in
\evensidemargin=0.25in
\setlength{\parindent}{0pt}

\begin{document}
\maketitle

\begin{enumerate}
    \item Done!
    \item Essentially, the planes which returned to England were those that suffered
          damage to noncritical areas; all those which received damage to other areas
          didn't come home. Thus, the areas which showed damage are noncritical, and
          the areas left undamaged on the returning planes should be reinforced.
    \item (a) is more probable, as for (b) to be true (a) must also be true. Thus
          there is no scenario where (b) occurs and (a) does not while there is a scenario
          where (a) occurs and (b) does not (that is, (b) occurs only if (a) does), so (b)
          is necessarily less likely than (a).
\end{enumerate}

\end{document}